\chapter{2023年HDU线性代数期中试题解析}
\section{填空题（每题3分，共18分）}
\vspace{-1em}
\begin{multicols}{2}
\begin{exercise}
    $
    \begin{vmatrix}
    5 & 2 & 0 & 0\\
    2 & 1 & 0 & 0\\
    0 & 0 & 8 & 3\\
    0 & 0 & 5 & 2
    \end{vmatrix}=
    $\fillans{$1$}.
\end{exercise}
\begin{solution}
$\begin{vmatrix}
    \boldsymbol{A} & \boldsymbol{O}\\
    \boldsymbol{O} & \boldsymbol{B}
\end{vmatrix}=\abs{\boldsymbol{AB}}=\abs{\boldsymbol{A}}\abs{\boldsymbol{B}}=1$.
\end{solution}
\end{multicols}

\begin{exercise}
    设矩阵$\boldsymbol{A}$满足$\boldsymbol{A}^2-\boldsymbol{A}-2\boldsymbol{E}=\boldsymbol{O}$，则$\ab(\boldsymbol{A}+2\boldsymbol{E})^{-1}$=\fillans{$\frac{\ab(3\boldsymbol{E}-\boldsymbol{A})}{4}$}.
\end{exercise}
\begin{solution}
    $\boldsymbol{A}^2-\boldsymbol{A}-2\boldsymbol{E}=\boldsymbol{O}$可分解为$\ab(\boldsymbol{A}+2\boldsymbol{E})\ab(3\boldsymbol{E}-\boldsymbol{A})=4\boldsymbol{E}$，由于零矩阵没有逆矩阵，故$\ab(\boldsymbol{A}+2\boldsymbol{E})^{-1}=\frac{\ab(3\boldsymbol{E}-\boldsymbol{A})}{4}$，
\end{solution}

\begin{exercise}
    已知$\boldsymbol{A}$为三阶方阵且$\abs{\boldsymbol{A}}=2$，则$\abs{\frac{1}{2}\boldsymbol{A}^*}=$\fillans{?}.
\end{exercise}
\begin{solution}
    $\abs{\frac{1}{2}\boldsymbol{A}^*}=\abs{\frac{1}{2}\abs{\boldsymbol{A}}\boldsymbol{A}^{-1}}=\abs{\boldsymbol{A}^{-1}}=\frac{1}{2}$.
\end{solution}

\begin{exercise}
$f\ab(x)=
\begin{vmatrix}
    5x & 1 & 2 & 3\\
    1 & x & 1 & 2\\
    1 & 2 & x & 3\\
    x & 1 & 2 & 2x
\end{vmatrix}
$的$x^3$项的系数是\fillans{$-3$}.
\end{exercise}
\begin{solution}
    $x^3$项为$x\cdot x\cdot x\cdot 3\cdot\ab(-1)^{\tau\ab(4231)}=-3x^3$
\end{solution}

\begin{exercise}
设$\boldsymbol{A}$是$4\times 3$矩阵，且$\mathrm{R}\ab(\boldsymbol{A})=2$，而$\boldsymbol{B}=
\begin{pmatrix}
9 & 0 & 0\\
9 & -7 & 0\\
8 & 7 & 8    
\end{pmatrix}$，则$\mathrm{R}\ab(\boldsymbol{AB})=$\fillans{$?$}.
\end{exercise}
\begin{solution}
    
\end{solution}

\begin{exercise}
    方阵$\boldsymbol{A}$为可逆矩阵的充要条件是\fillans{$\boldsymbol{A}$行等价于$\boldsymbol{E}$}.
\end{exercise}
\begin{solution}
    
\end{solution}

\section{选择题（每题3分，共18分）}
\vspace{-1em}
\begin{exercise}
    设$\boldsymbol{A}$是$n\times n$阶矩阵，且$\mathrm{R}\ab(\boldsymbol{A})<n$，则下面说法不正确的是\ans{}
    \begin{tasks}(2)
        \task $\abs{\boldsymbol{A}}=0$
        \task $\boldsymbol{A}$必有一个零子式
        \task 齐次线性方程组$\boldsymbol{AX}=0$仅有零解
        \task $\boldsymbol{A}$经初等行变换可化为单位矩阵
    \end{tasks}
\end{exercise}
\begin{solution}
    
\end{solution}

\begin{exercise}
    已知$\vv{\alpha}_1,\ \vv{\alpha}_2$为$2$维列向量，矩阵$\boldsymbol{A}=\ab(2\vv{\alpha}_1+\vv{\alpha}_2,\vv{\alpha}_1-\vv{\alpha}_2)$，$\boldsymbol{B}=\ab(\vv{\alpha}_1,\vv{\alpha}_2)$.若行列式$\abs{\boldsymbol{A}}=6$，则行列式$\abs{\boldsymbol{B}}=$\ans{A}
    \begin{tasks}(4)
        \task $-2$
        \task $2$
        \task $-3$
        \task $3$
    \end{tasks}
\end{exercise}
\begin{solution}
    
\end{solution}

\begin{exercise}
    设$\boldsymbol{A}$为$n$阶矩阵，对矩阵$\boldsymbol{A}$作若干次初等变换得到矩阵$\boldsymbol{B}$，那么必有\ans{}
    \begin{tasks}(4)
        \task $\abs{\boldsymbol{A}}=\abs{\boldsymbol{B}}$
        \task 若$\abs{\boldsymbol{A}}\neq 0$，则$\abs{\boldsymbol{B}}\neq 0$
        \task $\abs{\boldsymbol{A}}\neq\abs{\boldsymbol{B}}$
        \task 若$\abs{\boldsymbol{A}}>0$，则$\abs{\boldsymbol{B}}>0$
    \end{tasks}
\end{exercise}
\begin{solution}
    
\end{solution}

\begin{exercise}
    设$\boldsymbol{A}$为$m\times n$矩阵，$\boldsymbol{B}$为$\boldsymbol{n\times m}$矩阵，则线性方程组$\ab(\boldsymbol{AB})\boldsymbol{X}=0$\ans{}
    \begin{tasks}(2)
        \task 当$m>n$时仅有零解
        \task 当$m>n$时必有非零解
        \task 当$m<n$时仅有零解
        \task 当$m<n$时必有非零解
    \end{tasks}
\end{exercise}
\begin{solution}
    
\end{solution}

\begin{exercise}
    设$\boldsymbol{A},\ \boldsymbol{B},\ \boldsymbol{C}$是$n$阶方阵，且$\boldsymbol{ABC}=\boldsymbol{E}$，其中$\boldsymbol{E}$为单位矩阵，则\ans{C}
    \begin{tasks}(4)
        \task $\boldsymbol{ACB}=\boldsymbol{E}$
        \task $\boldsymbol{BAC}=\boldsymbol{E}$
        \task $\boldsymbol{BCA}=\boldsymbol{E}$
        \task $\boldsymbol{CBA}=\boldsymbol{E}$
    \end{tasks}
\end{exercise}
\begin{solution}
    $\epsilon_{\boldsymbol{A,B,C}}=1$
\end{solution}

\begin{exercise}
    设矩阵$\boldsymbol{A}=
    \begin{pmatrix}
        1 & x & x^2 & x^3\\
        1 & y & y^2 & y^3\\
        1 & z & z^2 & z^3
    \end{pmatrix}
    $，其中$x,\ y,\ z$互补相等，矩阵$\boldsymbol{A}$的秩等于\ans{}
    \begin{tasks}(4)
        \task $1$
        \task $2$
        \task $3$
        \task $4$
    \end{tasks}
\end{exercise}
\begin{solution}
    
\end{solution}

\section{计算题i（每题6分，共24分）}
\vspace{-1em}
\begin{exercise}
求四阶行列式$
\begin{vmatrix}
    2 & 1 & 4 & 1\\
    3 & -1 & 2 & 1\\
    1 & 2 & 3 & 2\\
    5 & 0 & 6 & 2
\end{vmatrix}
$
\end{exercise}
\begin{solution}
$
\begin{aligned}
    \begin{vmatrix}
        2 & 1 & 4 & 1\\
        3 & -1 & 2 & 1\\
        1 & 2 & 3 & 2\\
        5 & 0 & 6 & 2
    \end{vmatrix}\xlongequal{r_4-r_1-r_2}
    \begin{vmatrix}
        2 & 1 & 4 & 1\\
        3 & -1 & 2 & 1\\
        1 & 2 & 3 & 2\\
        0 & 0 & 0 & 0
    \end{vmatrix}=0
\end{aligned}
$
\end{solution}

\begin{exercise}
    试求矩阵$\boldsymbol{A}=
    \begin{pmatrix}
        2 & -1 & -1 & 1 & 2\\
        1 & 1 & -2 & 1 & 4\\
        4 & -6 & 2 & -2 & 4\\
        3 & 6 & -9 & 7 & 9
    \end{pmatrix}
    $的行最简形矩阵.
\end{exercise}
\begin{solution}
    $\begin{pmatrix}
        2 & -1 & -1 & 1 & 2\\
        1 & 1 & -2 & 1 & 4\\
        4 & -6 & 2 & -2 & 4\\
        3 & 6 & -9 & 7 & 9
    \end{pmatrix}
    \stackrel{r_2\leftrightarrow r_4}{\longrightarrow}
    \begin{pmatrix}
        2 & -1 & -1 & 1 & 2\\
        3 & 6 & -9 & 7 & 9\\
        4 & -6 & 2 & -2 & 4\\
        1 & 1 & -2 & 1 & 4
    \end{pmatrix}
    \stackrel{...}{\longrightarrow}\cdots
    $
\end{solution}

\begin{exercise}
    设四阶行列式$
    \begin{vmatrix}
        3 & -5 & 2 & 1\\
        1 & 1 & 0 & -5\\
        -1 & 3 & 1 & 3\\
        2 & -4 &-1 & -3
    \end{vmatrix}
    $，求$M_{11}+M_{21}+M_{31}+M_{41}$.
\end{exercise}
\begin{solution}
    $M_{11}+M_{21}+M_{31}+M_{41}=  
    \begin{vmatrix}
        1 & -5 & 2 & 1\\
        -1 & 1 & 0 & -5\\
        1 & 3 & 1 & 3\\
        -1 & -4 &-1 & -3
    \end{vmatrix}=0$
\end{solution}

\begin{exercise}
    已知$a$是常数，且矩阵$\boldsymbol{A}=
    \begin{pmatrix}
        1 & 2 & 2\\
        1 & 3 & 0\\
        2 & 7 & -2
    \end{pmatrix}
    $可经初等变换化为矩阵$\boldsymbol{B}=
    \begin{pmatrix}
        1 & a & 2\\
        0 &1 & 1\\
        -1 & 1 & 1
    \end{pmatrix}
    $，试求$a$的取值.
\end{exercise}
\begin{solution}
    
\end{solution}
\section{计算题ii（每题6分，共24分）}
\vspace{-1em}
\begin{exercise}
    设矩阵$\boldsymbol{A}=
    \begin{pmatrix}
        1 & -2 & 3 & -1 & 1\\
        3 & -1 & 5 & -3 & 2\\
        2 & 1 & 2 & -2 & 1
    \end{pmatrix}
    $，求矩阵$\boldsymbol{A}$的秩.
\end{exercise}
\begin{solution}
    $$\boldsymbol{A}=
    \begin{pmatrix}
        1 & -2 & 3 & -1 & 1\\
        3 & -1 & 5 & -3 & 2\\
        2 & 1 & 2 & -2 & 1
    \end{pmatrix}\longrightarrow
    \begin{pmatrix}
        1 & -2 & 3 & -1 & 1\\
        0 & 5 & -4 & 0 & -1\\
        0 & 5 & -4 & 0 & -1
    \end{pmatrix}\longrightarrow
    \begin{pmatrix}
        1 & -2 & 3 & -1 & 1\\
        0 & 5 & -4 & 0 & -1\\
        0 & 0 & 0 & 0 & 0
    \end{pmatrix}
    $$
    故$\mathrm{R}\ab(\boldsymbol{A})=2$.
\end{solution}

\begin{exercise}
    已知$\boldsymbol{A}=
    \begin{pmatrix}
        1 & 2 & 1\\
        3 & 4 & 2\\
        1 & 2 & 2
    \end{pmatrix}
    $，矩阵$\boldsymbol{X}$满足$\boldsymbol{AX}=\boldsymbol{A}^{\top}+\boldsymbol{X}$，求$\boldsymbol{X}$.
\end{exercise}
\begin{solution}
    $$
    \begin{aligned}
    \boldsymbol{X}&=\ab(\boldsymbol{A}-\boldsymbol{E})^{-1}\boldsymbol{A}^{\top}\\
    &=\begin{pmatrix}
        -1 & 0 & 1\\
        -1 & -1 & 3\\
        3 & 2 & -6
    \end{pmatrix}
    \begin{pmatrix}
        1 & 3 & 1\\
        2 & 4 & 2\\
        1 & 2 & 2
    \end{pmatrix}=
    \begin{pmatrix}
        0 & 0 & 1\\
        -1 & 0 & 3\\
        3 & 2 & -5
    \end{pmatrix}
    \end{aligned}
    $$
\end{solution}

\begin{exercise}
    设矩阵$\boldsymbol{A}=
    \begin{pmatrix}
        0 & 3 & 3\\
        1 & 1 & 0\\
        -1 & 2 & 3
    \end{pmatrix}
    $，且$\boldsymbol{AB}=\boldsymbol{A}+2\boldsymbol{B}$，求$\abs{\boldsymbol{B}}$.
\end{exercise}
\begin{solution}
    由矩阵方程
    $$\boldsymbol{AB}=\boldsymbol{A}+2\boldsymbol{B}$$
    解得
    $$\boldsymbol{B}=\ab(\boldsymbol{A}-2\boldsymbol{E})^{-1}\boldsymbol{A}$$
    $$\ab(\boldsymbol{A}-2\boldsymbol{E})^{-1}\boldsymbol{A}=\frac{1}{2}
    \begin{pmatrix}
        -1 & 3 & 3\\
        -1 & 1 & 3\\
        1 & 1 & -1
    \end{pmatrix}
    \begin{pmatrix}
        0 & 3 & 3\\
        1 & 1 & 0\\
        -1 & 2 & 3
    \end{pmatrix}
    =
    \begin{pmatrix}
        0 & 3 & 3\\
        -1 & 2 & 6\\
        1 & 1 & 0
    \end{pmatrix}
    $$
\end{solution}

\begin{exercise}
    已知矩阵$\boldsymbol{A}=
    \begin{pmatrix}
        2 & 1 & -3\\
        1 & 2 & -2\\
        -1 & 3 & 2
    \end{pmatrix}
    $，试求$\boldsymbol{A}$的逆矩阵.
\end{exercise}
\begin{solution}
    $$\ab(\boldsymbol{A}\vdots\boldsymbol{E})=
    \begin{pmatrix}
        2 & 1 & -3 &\vdots\ 1 & 0 & 0\\
        1 & 2 & -2 &\vdots\ 0 & 1 & 0\\
        -1 & 3 & 2 &\vdots\ 0 & 0 & 1
    \end{pmatrix}\longrightarrow
    \begin{pmatrix}
        1 & 0 & 0 &\vdots\ 2 & -11/5 & 4/5\\
        0 & 1 & 0 &\vdots\ 0 & 1/5 & 1/5\\
        0 & 0 & 1 &\vdots\ 1 & -7/5 & 3/5
    \end{pmatrix}=\ab(\boldsymbol{E}\vdots\boldsymbol{A}^{-1})
    $$
    $$\boldsymbol{A}^{-1}=
    \begin{pmatrix}
    2 & -11/5 & 4/5\\
    0 & 1/5 & 1/5\\
    1 & -7/5 & 3/5
    \end{pmatrix}
    $$
\end{solution}
\section{综合题（12分）}
\vspace{-1em}
\begin{exercise}
    已知线性方程组$
    \left\{
        \begin{aligned}
            &x_1+x_2+x_3=1\\
            &2x_1+5x_2+px_3=5\\
            &3x_1+px_2+3x_3=2
        \end{aligned}
    \right.
    $，问
    \vspace{1em}
    \begin{enumerate}
        \item 当参数$p$满足什么条件时，方程组有唯一解？无解？有无穷多解？
        \item 有无穷多解时，求方程组的通解.
    \end{enumerate}
\end{exercise}
\begin{solution}
    \begin{enumerate}
        \item 矩阵方程为
        $$
        \begin{pmatrix}
            1 & 1 & 1\\
            2 & 5 & p\\
            3 & p & 3
        \end{pmatrix}
        \begin{pmatrix}
            x_1\\x_2\\x_3
        \end{pmatrix}=
        \begin{pmatrix}
            1\\5\\2
        \end{pmatrix}
        $$
    \end{enumerate}
\end{solution}